\mxnewpage\section*{EV1-T2 -- DS commun}\label{EV1-T2}
\addcontentsline{toc}{section}{EV1-T2 -- DS commun}
\textbf{Sujet d’entraînement}. Chapitres 1 à 4.\par Durée : 2 h. Sans calculatrice. Note sur 20. 14 novembre 2026.
\subsection*{Partie 1 -- Automatismes (6 points)}
Douze questions, 0{,}5 point chacune. Une seule réponse correcte ; aucune justification demandée. Une réponse incorrecte ou absente rapporte 0, sans retrait de points.
\begin{enumerate}
\item Calculer $4+2\times3$.\par A. $18$\quad B. $9$\quad C. $13$\quad D. $10$
\item Calculer $\frac12+\frac13$.\par A. $\frac{5}{3}$\quad B. $\frac{2}{5}$\quad C. $\frac{5}{6}$\quad D. $\frac{1}{6}$
\item Calculer $\dfrac{2^4}{2^{2}}$.\par A. $2$\quad B. $4$\quad C. $1$\quad D. $9$
\item Une valeur de 50 augmente de 20 \%. Valeur finale ?\par A. $70$\quad B. $40$\quad C. $10$\quad D. $60$
\item Une valeur de 80 est obtenue après une baisse de 20 \%. Valeur initiale ?\par A. $100$\quad B. $96$\quad C. $64$\quad D. $60$
\item Hausse de 25 \% puis baisse de 20 \% : coefficient global ?\par A. $1$\quad B. $\frac{21}{20}$\quad C. $\frac{4}{5}$\quad D. $\frac{9}{10}$
\item La solution de $2x+4=0$ vaut\par A. $2$\quad B. $8$\quad C. $-2$\quad D. $-8$
\item La pente entre $(0;4)$ et $(2;10)$ vaut\par A. $3$\quad B. $-3$\quad C. $6$\quad D. $2$
\item La moyenne de $4, 6, 8$ vaut\par A. $18$\quad B. $6$\quad C. $7$\quad D. $4$
\item Un secteur de $90^\circ$ représente quelle fraction du total ?\par A. $\frac{1}{3}$\quad B. $\frac{3}{4}$\quad C. $\frac{1}{4}$\quad D. $\frac{1}{2}$
\item Une probabilité vaut $\dfrac{2}{5}$. Son complément vaut\par A. $\frac{5}{2}$\quad B. $\frac{2}{5}$\quad C. $\frac{1}{5}$\quad D. $\frac{3}{5}$
\item Pour $f(x)=x^2+4$, $f(-2)$ vaut\par A. $6$\quad B. $16$\quad C. $8$\quad D. $0$
\end{enumerate}
\mxnewpage\subsection*{Partie 2 -- Exercices indépendants (14 points)}
Toute démarche doit être justifiée. Les résultats exacts sont attendus sauf indication contraire.
\mxneed{12}
\subsection*{Exercice 1 -- Deux procédés de croissance (5 points)}
On compare deux modèles pour la production d’un atelier. Le jour 0, chacun prévoit 30 objets. Le modèle A prévoit 10 objets supplémentaires chaque jour ; le modèle B prévoit un doublement quotidien. Les productions du jour $n$ sont donc $a_n=30+10n$ et $b_n=30\times2^n$, pour $n\in\N$. Un cumul désigne ici la somme des productions des jours indiqués.
\begin{enumerate}
\item Donner la nature de chaque suite et calculer les productions du jour 2. \hfill(1 pt)
\item Calculer les productions cumulées des jours 0 à 3 dans chacun des modèles. \hfill(2 pts)
\item Vérifier que le modèle B dépasse le modèle A au jour 1, puis expliquer pourquoi cet écart reste positif à tous les jours suivants en comparant les accroissements quotidiens. \hfill(2 pts)
\end{enumerate}
\mxneed{12}
\subsection*{Exercice 2 -- Des racines imposées (5 points)}
On cherche une fonction polynôme $f$ de degré 2, définie sur $\R$. Sa courbe coupe l’axe des abscisses aux abscisses $-1$ et 5, et elle passe par le point $(0;10)$. Ces informations doivent permettre de déterminer la fonction.
\begin{enumerate}
\item Donner sa forme factorisée. \hfill(1 pt)
\item Déterminer son signe puis résoudre $f(x)\leq0$. \hfill(2 pts)
\item Donner sa forme canonique et son maximum. \hfill(2 pts)
\end{enumerate}
\mxneed{12}
\subsection*{Exercice 3 -- Une approximation contrôlée (4 points)}
La fonction $f$ est définie sur $\R$ par $f(x)=x^2+1$. On souhaite approcher $f(3{,}02)$ à l’aide de la tangente à sa courbe au point d’abscisse 3. Les calculs sont effectués sans calculatrice.
\begin{enumerate}
\item Pour $h\ne0$, calculer et simplifier le taux de variation entre 3 et $3+h$. \hfill(1 pt)
\item En déduire $f\prime(3)$, puis l’équation de la tangente au point d’abscisse 3. \hfill(1 pt)
\item Utiliser la tangente pour approcher $f(3{,}02)$. Calculer ensuite l’erreur « valeur exacte moins approximation ». \hfill(2 pts)
\end{enumerate}
